\section{Distillation Results：SFT 是地基，DPO 是增量}

在 human-data ablation 之后，课程回到 Zephyr。这里需要把 leaderboard claim 与 training ablation 分开：前者说明特定协议上的竞争力，后者才更直接回答 dSFT、chosen-only SFT 与 dDPO 各自贡献。

\subsection{Lecture-time Leaderboard Evidence}

Zephyr-7B 基于 Mistral-7B，使用 UltraChat 做 dSFT，再使用 UltraFeedback chosen/rejected 做 dDPO。课堂展示它在当时 MT-Bench 与 AlpacaEval 上接近或超过更大模型；这些结果依赖 GPT-4 family judges，也正是下一章要审计的对象。

\lecturefigure{slide-56.jpg}{Distillation results 的章节转场}{官方幻灯片 p.56}

\begin{knowledgebox}{读图：从 human-curated ablation 切换到 AI-distilled recipe}
这一页虽然信息简洁，却标志证据来源变化：后续模型不再主要学习 vendor-written demonstrations，而是学习 AI-generated dialogues/preferences。比较时必须把 provenance 一起带上。
\end{knowledgebox}

\lecturefigure{slide-57.jpg}{Zephyr-7B 在 MT-Bench 与 AlpacaEval 的课堂结果}{官方幻灯片 p.57}

\begin{knowledgebox}{读图：红框说明协议内竞争力，不是全维度胜利}
表中同时包含参数规模、MT-Bench 与 AlpacaEval。先比较同规模 open models，再看与 proprietary systems 的 judge score；图没有直接测安全、事实核验、真实用户 retention 或分布外 robustness。
\end{knowledgebox}

\subsection{Ablation：DPO-only、SFT-only 与 SFT+DPO}

更有机制信息的是 ablation。直接在 base model 上做 DPO 表现很差；只做 dSFT 已获得大部分收益；在 dSFT 之后加 dDPO 再提高一部分。把 chosen responses 当普通 SFT targets 也不能完全替代 preference objective。

\lecturefigure{slide-58.jpg}{Zephyr 的 dSFT/dDPO ablation}{官方幻灯片 p.58}

\begin{knowledgebox}{读图：先找必要条件，再比较增量}
DPO-only 行说明 preference optimization 需要已具备 instruction-following 的 policy；SFT-only 行约获得最终表现的 80--90\%；SFT+DPO 最好但增量较小。该表支持“DPO 是 refinement”，不支持“DPO 无用”或“任何 SFT 都足够”。
\end{knowledgebox}

\teachervoice{讲者明确说 DPO without SFT ``doesn't work at all''，SFT 已带来 80--90\% 的整体表现，而 DPO 提供剩余增益。Chosen-only SFT 也不等价于利用 pairwise preference。}

\subsection{本章小结}

Zephyr 证明 AI-generated data 能支持强 open model，但也把 teacher/judge bias 带入 pipeline。Ablation 的持久结论是：instruction-following foundation 先由 SFT 建立，preference optimization 再做相对精炼。

